% Generated by tools/edn_latex.py from reports/edn.md.
% Do not edit: change reports/edn.md and run `make edn`.
\ednentry{
  id = {EDN-10},
  title = {Availability target replaced by recovery time plus a presentation-window commitment},
  date = {2026-09-22},
  milestone = {M6: Monitoring},
  topic = {Resource Monitoring, ML System Design},
  participants = {@lukas2510},
  decision = {NFR-05 no longer commits to ``99\,\% uptime in the weeks before each presentation.'' It now commits to: (1) automatic recovery, the API reaches a healthy \texttt{/\allowbreak{}health} within 30 s of a crash or VM reboot, verified by a chaos test; (2) zero unplanned downtime during each live presentation and its warm-up, checked manually right before presenting; (3) outside those windows, uptime on the single, unattended, no-SLA VM is monitored and reported via Better Uptime and Grafana, not contractually targeted. Better Uptime is added to the architecture, the M6 tool list and the reference links, none of which named it before.},
  rationale = {\emph{Alternatives considered:}
\begin{itemize}[leftmargin=*, topsep=2pt]
\item \textbf{Option A: keep 99\,\% uptime over ``the weeks before each presentation.''}
\newline Pros: a single, simple number; easy to write into the report.
\newline Cons: nothing backs it: one VM, no redundancy, no SLA from FIB Virtech, no on-call team; the only listed check (Prometheus \texttt{up}) can't see the host itself going down; the measurement window is unbounded and unverifiable after the fact, so the number could not actually be defended if a teacher asked how it was checked.
\item \textbf{Option B: lower the aggregate target (e.g. 95\,\%) instead of dropping it.}
\newline Pros: keeps a single reportable number; slightly more honest than 99\,\%.
\newline Cons: still as unfalsifiable as 99\,\% without redundancy or an on-call team; still hides that the requirement was actually about two specific presentation days, not a general SLA the team can neither guarantee nor is expected to for a course project.
\item \textbf{Option C (chosen): split into recovery time, a narrow presentation-window commitment, and a monitored-not-targeted general period.}
\newline Pros: the recovery-time claim is testable today with a chaos test, not only after weeks of passive observation; the presentation-window claim is the one that actually matters for grading and is narrow enough for the team to actively watch; the monitored-general-period framing is honest about the single-VM, no-redundancy, no-SLA deployment while still giving the M6 monitoring practice something to show (Better Uptime + Grafana), which is what the rubric actually grades, not a specific SLA number.
\newline Cons: three clauses instead of one number, slightly longer to state.
\end{itemize}
\par
\emph{Why:} M6's Resource Monitoring criterion asks whether the monitoring anticipates real-world challenges and gives actionable insight, not whether a specific SLA percentage was hit. A number the team has no redundancy or on-call capacity to guarantee is a liability if questioned in the presentation, not evidence of the practice. Splitting the requirement matches what is actually controllable (recovery time, presentation-day behaviour) from what can only be observed and reported (general uptime on a single, unattended, no-SLA VM), and it adds Better Uptime, a tool the course's own M6 schedule already names but the docs previously omitted.},
  ai = {Information seeking, Alternative generation, Alternative assessment, Recommendation},
  response = {Accepted},
  assessment = {During the PR \#19 review, AI had already flagged that Prometheus's \texttt{up} metric cannot observe the host going down and that Better Uptime was missing from the docs despite being a named M6 course tool. Asked directly what availability target would actually be feasible given the single-VM, no-redundancy deployment, AI proposed the three-part split (recovery time, presentation window, monitored general period) with rationale grounded in the M6 rubric wording and the deployment's real constraints, and named option B as a weaker alternative. Lukas reviewed it and confirmed the split.},
  aievidence = {Claude Code session on 2026-09-22 while reviewing PR \#19: Lukas said ``i dont like that [99\% uptime]. what NFR that goes in the uptime direction is feasible with our system and constraints?''; AI proposed the three-part split and the addition of Better Uptime; Lukas replied ``ok do that.''},
  evidence = {\href{https://github.com/mlops-2627q1-mds-upc/MLOPS_Recommenditos/blob/main/docs/docs/requirements.md}{requirements} NFR-05; \href{https://github.com/mlops-2627q1-mds-upc/MLOPS_Recommenditos/blob/main/docs/docs/specification.md}{specification} NFR-05; \href{https://github.com/mlops-2627q1-mds-upc/MLOPS_Recommenditos/blob/main/docs/docs/project-brief.md}{project brief} section 5 (target architecture) and section 7 (milestones); \href{https://github.com/mlops-2627q1-mds-upc/MLOPS_Recommenditos/blob/main/references/MLOps-lab.md}{MLOps-lab.md} session 11; \href{https://github.com/mlops-2627q1-mds-upc/MLOPS_Recommenditos/pull/19}{PR \#19}.},
}
